% Apunts de teoria en PDF · 4t ESO · Unitat 10: Estadística bidimensional
% Mateix estil que els apunts de Batxillerat (franja i caixes de teoria en blau, exemples en roig), però més visuals,
% curts i amb dibuixos, per a l'alumnat de 3r ESO. Es compila amb pdflatex dues vegades (per als enllaços de l'índex),
% amb logo-fentmates.pdf (_eines/logo/pdf/horitzontal-clar.pdf) al costat.
\documentclass[11pt,a4paper]{article}
\usepackage[utf8]{inputenc}
\usepackage[T1]{fontenc}
\usepackage[catalan]{babel}
\usepackage[sfdefault,book,lining]{FiraSans}
\usepackage[italic,noplusnominus,noequal]{mathastext}
\linespread{1.18}
\usepackage[a4paper,top=2.4cm,bottom=2.2cm,left=1.6cm,right=1.6cm,headheight=1.1cm,headsep=0.5cm,footskip=1cm]{geometry}
\usepackage{amsmath,amssymb,cancel}
\usepackage{icomma}
\usepackage{graphicx}
\usepackage[table]{xcolor}
\usepackage{tikz}\usetikzlibrary{calc,positioning,arrows.meta,decorations.pathreplacing,shapes.geometric}
\usepackage{fontawesome5}
\usepackage{ccicons}
\usepackage{fancyhdr}
\usepackage{enumitem}
\usepackage{array}
\usepackage{needspace}
\usepackage[most]{tcolorbox}
\usepackage{titlesec}
\usepackage[hidelinks]{hyperref}
\hypersetup{pdftitle={Apunts de teoria. Unitat 10: Estadística bidimensional. 4t ESO},pdfsubject={Matemàtiques A 4t ESO},pdfkeywords={variable bidimensional, taula de doble entrada, núvol de punts, centre de gravetat, correlació, coeficient de correlació, covariància, recta de regressió, estimacions, interpolació, extrapolació}}

% colors de la web (bloc d'Àlgebra: blau)
\definecolor{fmnavy}{HTML}{1B2748}
\definecolor{fmblue}{HTML}{3D6AA0}
\definecolor{fmbluel}{HTML}{4E81BD}
\definecolor{fmyellow}{HTML}{FFC94D}
\definecolor{fmgold}{HTML}{E9A919}
\definecolor{fmaccent}{HTML}{FF7A50}
\definecolor{fmdef}{HTML}{1D4F73}\definecolor{fmdefbg}{HTML}{EEF4F9}
\definecolor{fmex}{HTML}{A3203A}\definecolor{fmexbg}{HTML}{FBF0F2}
\definecolor{fmsf}{HTML}{C2410C}
\definecolor{fmok}{HTML}{15803D}
\definecolor{fmmuted}{HTML}{5D6685}
\definecolor{fmsurface}{HTML}{EEF1F8}
\definecolor{fmA}{HTML}{2F6FD6}\definecolor{fmB}{HTML}{EA6A10}\definecolor{fmR}{HTML}{E11D48}\definecolor{fmK}{HTML}{16A34A}\definecolor{fmV}{HTML}{7C3AED}

% lletres
\newcommand{\comic}{\fontfamily{ComicNeue-TLF}\selectfont}
\newcommand{\heavy}{\fontseries{eb}\selectfont}
\newcommand{\semi}{\fontseries{sb}\selectfont}
\DeclareUnicodeCharacter{2212}{-}

\setlength{\parindent}{0pt}
\setlength{\parskip}{7pt}
\setlength{\lineskip}{4pt}\setlength{\lineskiplimit}{4pt}
\setlength{\emergencystretch}{3em}
\setlist{itemsep=3pt,topsep=3pt,leftmargin=1.4em}
\setcounter{secnumdepth}{1}
\renewcommand{\thesection}{10.\arabic{section}}
\titleformat{\section}[hang]{\color{fmnavy}\fontsize{23}{27}\heavy\selectfont}{\tcbox[on line,colback=fmblue,colframe=fmblue,arc=3mm,boxsep=2pt,left=6pt,right=6pt,top=3pt,bottom=3pt]{\color{white}\thesection}}{0.45em}{}[\vspace{5pt}{\color{fmblue}\titlerule[2pt]}]
\titlespacing*{\section}{0pt}{4pt}{6pt}
\titleformat{\subsection}{\color{fmblue}\Large\semi}{}{0pt}{\faCaretRight\ }
\titlespacing*{\subsection}{0pt}{8pt}{2pt}

% frase de cada apartat, «a mà»
\newcommand{\frase}[1]{\par\vspace{-2pt}{\comic\Large\color{fmbluel}\faCommentDots[regular]\ \ #1\par}\vspace{2pt}}
\newcommand{\seccio}[3]{\Needspace{16\baselineskip}\section{#2}\hypertarget{#1}{}\frase{#3}}
% notes a mà dins dels exemples
\newcommand{\nota}[1]{{\comic\small\color{fmex!85!black}#1}}
\newcommand{\notab}[1]{{\comic\small\color{fmblue}#1}}
% bé / malament
\newcommand{\be}[1]{{\color{fmok}\faCheckCircle}\ \ensuremath{#1}}
\newcommand{\mal}[1]{{\color{fmR}\faTimesCircle}\ {\color{fmR}\ensuremath{#1}}}
% pas numerat
\newcommand{\pas}[1]{\tikz[baseline=(p.base)]\node[circle,fill=fmblue,text=white,inner sep=1.2pt,minimum size=1.25em,font=\small\heavy](p){#1};}
% etiqueta de color
\newcommand{\chip}[2]{\tcbox[on line,colback=#1!14,colframe=#1!60,arc=2mm,boxsep=1pt,left=5pt,right=5pt,top=1pt,bottom=1pt]{#2}}
\newcommand{\fmkey}[1]{\colorbox{fmyellow!55}{\strut #1}}
\newcommand{\cA}[1]{{\color{fmA}#1}}\newcommand{\cB}[1]{{\color{fmB}#1}}\newcommand{\cK}[1]{{\color{fmK}#1}}\newcommand{\cR}[1]{{\color{fmR}#1}}\newcommand{\cV}[1]{{\color{fmV}#1}}

% caixes
\tcbset{fmbase/.style={enhanced,arc=3mm,boxrule=1.2pt,left=4mm,right=4mm,top=2mm,bottom=2.5mm,before skip=8pt,after skip=8pt,fonttitle=\heavy\large,before upper=\raggedright}}
\newtcolorbox{idea}[1][Idea clau]{fmbase,colback=fmdefbg,colframe=fmdef,coltitle=white,title={\faKey\ \ #1}}
\newtcolorbox{exemple}[1][Exemple]{fmbase,colback=fmexbg,colframe=fmex,colbacktitle=fmexbg,coltitle=fmex,title={\faPencil*\ \ #1},titlerule=0.5pt,titlerule style={fmex!40,dashed}}
\newtcolorbox{compte}[1][Compte!]{fmbase,colback=fmsf!6,colframe=fmsf,coltitle=white,fonttitle=\comic\bfseries\Large,title={\faExclamationTriangle\ \ #1}}
\newtcolorbox{truc}[1][Truc]{fmbase,colback=fmyellow!20,colframe=fmgold,colbacktitle=fmyellow,coltitle=fmnavy,fonttitle=\comic\bfseries\Large,title={\faLightbulb[regular]\ \ #1}}

% dibuix de la targeta de la unitat, en blanc
\newcommand{\fmcard}{\begin{tikzpicture}[x=0.062cm,y=-0.062cm,white,line width=3pt,line cap=round]
  \begin{scope}[transparency group,opacity=0.15]
  \node[anchor=base west,inner sep=0pt,font=\fontsize{28}{28}\selectfont] at (14,44) {$r\approx0{,}9$};
  \draw[line width=2pt] (120,110) -- (235,110) (125,115) -- (125,15);
  \foreach \p in {(140,95),(152,88),(160,92),(172,75),(180,70),(192,62),(200,50),(212,45),(222,32)} \fill \p circle (5);
  \draw[line width=2.4pt] (135,100) -- (230,28);
  \end{scope}
\end{tikzpicture}}
\newcommand{\fmbrand}{\raisebox{-0.25\height}{\includegraphics[height=0.5cm]{logo-fentmates.pdf}}}

% regla de Ruffini: #1 = a, #2 = coeficients, #3 = fila del mig (sense el primer), #4 = fila de baix, #5 = nombre de coeficients
\newcommand{\ruffini}[5]{\begin{tikzpicture}[x=1.25cm,y=0.85cm,every node/.style={font=\large}]
  \foreach \v [count=\i] in {#2} \node (t\i) at (\i,0) {$\v$};
  \foreach \v [count=\i] in {#3} {\pgfmathtruncatemacro\j{\i+1}\node (m\j) at (\j,-1) {$\v$};}
  \foreach \v [count=\i] in {#4} \node (b\i) at (\i,-2.1) {$\v$};
  \node[text=fmB] at (0,-1) {$#1$};
  \draw[line width=1.1pt,fmnavy] (0.45,0.5)--(0.45,-1.5) (0.45,-1.5)--(#5+0.5,-1.5);
  \node[draw=fmR,line width=1.3pt,rounded corners=2mm,minimum size=0.75cm] at (#5,-2.1) {};
  \draw[-{Latex[length=2.2mm]},fmK,line width=1.1pt] (t1.south) -- (b1.north);
  \pgfmathtruncatemacro\last{#5-1}
  \foreach \i in {1,...,\last} {\pgfmathtruncatemacro\j{\i+1}
    \draw[-{Latex[length=2.2mm]},fmB,line width=1.1pt,shorten >=1pt] (b\i.north east) -- (m\j.south west);
    \draw[-{Latex[length=2.2mm]},fmA!70,line width=1pt,dashed,shorten >=2pt] ([xshift=8pt]t\j.south) -- ([xshift=8pt]b\j.north);}
\end{tikzpicture}}

% capçalera i peu
\pagestyle{fancy}\fancyhf{}
\fancyhead[L]{\small\fmbrand}
\fancyhead[R]{\small\color{fmmuted}4t ESO · Unitat 10: Estadística bidimensional}
\fancyfoot[L]{\scriptsize\color{fmmuted}\href{https://creativecommons.org/licenses/by-nc-sa/4.0/deed.ca}{\ccbyncsa}\ \ FentMates · fentmates.org · CC BY-NC-SA 4.0}
\fancyfoot[R]{\small\color{fmmuted}\thepage}
\renewcommand{\headrulewidth}{0.4pt}\renewcommand{\headrule}{\color{fmmuted!40}\hrule height \headrulewidth}
\fancypagestyle{primera}{\fancyhf{}\renewcommand{\headrulewidth}{0pt}
  \fancyfoot[L]{\scriptsize\color{fmmuted}\href{https://creativecommons.org/licenses/by-nc-sa/4.0/deed.ca}{\ccbyncsa}\ \ FentMates · fentmates.org · CC BY-NC-SA 4.0}
  \fancyfoot[R]{\small\color{fmmuted}\thepage}}

% targeta de l'índex
\newcommand{\itile}[4]{\begin{tcolorbox}[enhanced,colback=white,colframe=fmblue!35,arc=3mm,boxrule=1pt,boxsep=0pt,left=6pt,right=6pt,top=5pt,bottom=5pt,before skip=0pt,after skip=7pt]
  \hyperlink{#1}{\begin{minipage}[c]{0.9cm}\centering{\color{fmblue}\LARGE #3}\end{minipage}\hspace{6pt}%
  \begin{minipage}[c]{0.78\linewidth}{\color{fmblue}\heavy #2}\\[-2pt]{\color{fmnavy}\semi\large #4}\end{minipage}}\end{tcolorbox}}

% taula amb vores arrodonides: el contingut es retalla amb la forma (sense @{} a les vores ni entre columnes)
\newsavebox\tabclipbox
\newcommand{\taulaclip}[3][3mm]{\sbox\tabclipbox{#3}\begin{tikzpicture}[baseline=(tcn.base)]
  \begin{scope}\clip[rounded corners=#1] (0,-\dp\tabclipbox) rectangle (\wd\tabclipbox,\ht\tabclipbox);
  \node[inner sep=0pt,outer sep=0pt,anchor=base west] (tcn) at (0,0) {\usebox\tabclipbox};\end{scope}
  \draw[#2,rounded corners=#1] (0,-\dp\tabclipbox) rectangle (\wd\tabclipbox,\ht\tabclipbox);
\end{tikzpicture}}
\newcommand{\taularod}[3][fmblue]{\taulaclip{#1,line width=1.2pt}{\renewcommand{\arraystretch}{1.5}\begin{tabular}{#2}#3\end{tabular}}}
% balança
\newcommand{\balanca}[3]{\begin{tikzpicture}[scale=#3,transform shape,every node/.style={anchor=center}]
  \fill[fmnavy] (-0.45,-1.9) rectangle (0.45,-1.75);
  \fill[fmnavy!80] (-0.12,-1.75) rectangle (0.12,0);
  \fill[fmgold] (0,0.05) circle (0.13);
  \draw[fmnavy,line width=2.4pt,line cap=round] (-2,0) -- (2,0);
  \foreach \s in {-1,1} {
    \draw[fmmuted,line width=0.8pt] (\s*2,0) -- (\s*2-0.75,-1.05) (\s*2,0) -- (\s*2+0.75,-1.05);
    \fill[fmyellow!70,draw=fmgold,line width=1pt] (\s*2-0.95,-1.05) arc[start angle=180,end angle=360,x radius=0.95,y radius=0.32] -- cycle;}
  \node[font=\Large,fill=fmA!15,draw=fmA,rounded corners=1.5mm,inner sep=3pt] at (-2,-0.68) {$#1$};
  \node[font=\Large,fill=fmB!15,draw=fmB,rounded corners=1.5mm,inner sep=3pt] at (2,-0.68) {$#2$};
\end{tikzpicture}}
% eixos amb quadrícula (dins d'un tikzpicture): xmin, xmax, ymin, ymax
\newcommand{\eixos}[4]{\draw[fmmuted!22,line width=0.4pt] (#1,#3) grid (#2,#4);
  \draw[-{Latex[length=2mm]},fmmuted,line width=0.8pt] (#1-0.2,0)--(#2+0.4,0) node[right,font=\small]{$x$};
  \draw[-{Latex[length=2mm]},fmmuted,line width=0.8pt] (0,#3-0.2)--(0,#4+0.4) node[above,font=\small]{$y$};}
% taula de traduccions (diccionari) amb les vores arrodonides
\newcommand{\dic}[2]{#1 & {\large$#2$}\\}
\newcommand{\dicratio}{0.62}
\newcommand{\dichead}{\begin{tabular}{@{}>{\raggedright}p{\dimexpr\dicratio\linewidth-12pt}>{\centering\arraybackslash}p{\dimexpr\linewidth-\dicratio\linewidth-12pt}@{}}\faComment[regular]\ Ho diem així\dots & \dots\ així\end{tabular}}
\newenvironment{dicbox}{\begin{tcolorbox}[enhanced,arc=3mm,boxrule=1.2pt,colframe=fmblue,colback=white,coltitle=white,fonttitle=\heavy,
  title={\dichead},left=0pt,right=0pt,top=0pt,bottom=1.5mm,boxsep=0pt,lefttitle=6pt,righttitle=6pt,before skip=0pt,after skip=0pt]\rowcolors{1}{white}{fmsurface}%
  \begin{tabular}{@{\hspace{6pt}}>{\raggedright}p{\dimexpr\dicratio\linewidth-12pt}>{\centering\arraybackslash}p{\dimexpr\linewidth-\dicratio\linewidth-12pt}@{\hspace{6pt}}}}
  {\end{tabular}\end{tcolorbox}}
% resum
\newcommand{\rtile}[3]{\begin{tcolorbox}[enhanced,colback=white,colframe=fmblue!45,arc=3mm,boxrule=1pt,left=3mm,right=3mm,top=2.5mm,bottom=2.5mm,fontupper=\large,
  title={\tcbox[on line,colback=fmblue,colframe=fmblue,arc=2mm,boxsep=0pt,left=4pt,right=4pt,top=1pt,bottom=1pt]{\color{white}\small\heavy #1}\ \ #2},coltitle=fmnavy,colbacktitle=fmdefbg,fonttitle=\semi,toptitle=0.5mm,bottomtitle=0.5mm]#3\end{tcolorbox}}
% portada: #1 subtítol a mà
\newcommand{\portada}[1]{\thispagestyle{primera}
\begin{tikzpicture}[remember picture,overlay]
  \fill[fmblue] (current page.north west) rectangle ([yshift=-9.4cm]current page.north east);
  \begin{scope}
    \clip (current page.north west) rectangle ([yshift=-9.4cm]current page.north east);
    \node[anchor=north west,inner sep=0pt] at ([xshift=6.6cm,yshift=-1.3cm]current page.north west) {\fmcard};
  \end{scope}
  \node[anchor=north west,inner sep=0pt] at ([xshift=1.6cm,yshift=-1.1cm]current page.north west) {\tcbox[colback=white,colframe=white,arc=2.5mm,boxsep=0pt,left=7pt,right=7pt,top=5pt,bottom=5pt]{\includegraphics[height=0.75cm]{logo-fentmates.pdf}}};
\end{tikzpicture}
\vspace*{0.55cm}
\tcbox[colback=white!15!fmblue,colframe=white!15!fmblue,arc=2.5mm,boxsep=0pt,left=6pt,right=6pt,top=3pt,bottom=3pt]{\color{white}\small\bfseries Apunts de teoria · 4t ESO · Matemàtiques A}

\vspace{2pt}
{\raggedright\color{white}\fontsize{30}{36}\heavy\selectfont{\color{fmyellow}Unitat 10}\\[2pt]Estadística bidimensional\par}
\vspace{8pt}
{\color{white}\comic\Large #1\par}
\vspace*{\dimexpr 7.4cm-\pagetotal\relax}
}
\newcommand{\llegenda}{\par\vspace{16pt}
{\color{fmnavy}\heavy\Large\faPalette\ \ Com es llegeixen aquests apunts}\par\smallskip
\begin{tcbraster}[raster columns=4,raster equal height,raster column skip=8pt,raster before skip=12pt]
\tcbox[colback=fmdefbg,colframe=fmdef,arc=3mm,boxrule=1.2pt,halign=center,width=\linewidth]{{\color{fmdef}\Large\faKey}\\[2pt]{\color{fmdef}\heavy Idea clau}\\[-1pt]{\small el que cal saber}}
\tcbox[colback=fmexbg,colframe=fmex,arc=3mm,boxrule=1.2pt,halign=center,width=\linewidth]{{\color{fmex}\Large\faPencil*}\\[2pt]{\color{fmex}\heavy Exemple}\\[-1pt]{\small pas a pas}}
\tcbox[colback=fmsf!6,colframe=fmsf,arc=3mm,boxrule=1.2pt,halign=center,width=\linewidth]{{\color{fmsf}\Large\faExclamationTriangle}\\[2pt]{\color{fmsf}\heavy Compte!}\\[-1pt]{\small errors típics}}
\tcbox[colback=fmyellow!20,colframe=fmgold,arc=3mm,boxrule=1.2pt,halign=center,width=\linewidth]{{\color{fmgold!80!black}\Large\faLightbulb[regular]}\\[2pt]{\color{fmnavy}\heavy Truc}\\[-1pt]{\small per a recordar-ho}}
\end{tcbraster}}
\newcommand{\resumtitol}{\clearpage
\hypertarget{resum}{}
{\color{fmnavy}\fontsize{23}{27}\heavy\selectfont\faClipboardList\ \ En resum}\par\vspace{5pt}{\color{fmblue}\titlerule[2pt]}
\frase{Tota la unitat en una pàgina.}}
% els 5 passos per a resoldre problemes: #2 = text del pas 2
\newcommand{\passosproblema}[2]{\begin{center}
\begin{tikzpicture}[s/.style={draw=fmblue,fill=fmdefbg,rounded corners=3mm,line width=1.2pt,minimum width=2.75cm,minimum height=2.45cm,align=center,font=\small},
  ic/.style={font=\LARGE,text=fmblue}]
  \foreach \i/\ic/\t in {1/\faBookOpen/{\textbf{Llegim}\\i entenem},2/\faQuestionCircle/{#1},3/\faEquals/{#2},4/\faCogs/{\textbf{Resolem}},5/\faCheckDouble/{\textbf{Comprovem}\\i responem}} {
    \node[s] (s\i) at (3.3*\i,0) {};
    \node[ic] at ([yshift=-0.5cm]s\i.north) {\ic};
    \node[align=center,font=\small,anchor=north] at ([yshift=-0.95cm]s\i.north) {\t};
    \node[circle,fill=fmblue,text=white,font=\small\heavy,inner sep=1.5pt,minimum size=1.3em] at (s\i.north west) {\i};}
  \foreach \i [evaluate=\i as \j using int(\i+1)] in {1,...,4} \draw[-{Latex[length=3mm]},fmmuted,line width=1.5pt] (s\i) -- (s\j);
\end{tikzpicture}
\end{center}}

\newcommand{\lo}{\mathopen{]}}\newcommand{\ro}{\mathclose{[}}
% punts ple i buit, nombres dels eixos i taula de valors
\newcommand{\pple}[2][fmblue]{\fill[#1] #2 circle (3.6pt);}
\newcommand{\nums}[4]{\foreach \k in {#1,...,#2} {\ifnum\k=0\else\node[below,font=\scriptsize,text=fmmuted] at (\k,0) {$\k$};\fi}
  \foreach \k in {#3,...,#4} {\ifnum\k=0\else\node[left,font=\scriptsize,text=fmmuted] at (0,\k) {$\k$};\fi}}
\newcommand{\tv}[3][fmblue]{\taulaclip{#1,line width=1.2pt}{\renewcommand{\arraystretch}{1.45}%
  \begin{tabular}{>{\columncolor{#1}\color{white}\heavy}c*{12}{w{c}{0.8cm}}}#2\\\rowcolor{#1!8}\cellcolor{#1}#3\end{tabular}}}
% etiqueta d'una recta o corba
\newcommand{\etq}[3]{\node[fill=white,inner sep=1.5pt,text=#1,font=\small] at #2 {#3};}

% interval sobre la recta: #1 esquerre, #2 dret, #3 tipus esquerre (o/c/a), #4 tipus dret (o/c/a), #5 color
\def\intvdot#1#2#3{\ifx#2o\filldraw[fill=white,draw=#3,line width=1.3pt] (#1,0) circle (3.3pt);\fi\ifx#2c\fill[#3] (#1,0) circle (3.5pt);\fi}
\newcommand{\intv}[5]{\begin{tikzpicture}[x=0.4cm,baseline=-3pt]
  \draw[fmmuted,line width=0.8pt] (-6,0)--(7,0);
  \foreach \k in {-5,...,6} \draw[fmmuted,line width=0.6pt] (\k,-0.08)--(\k,0.08);
  \ifx#3a\def\xa{-6}\else\def\xa{#1}\fi
  \ifx#4a\def\xb{7}\else\def\xb{#2}\fi
  \ifx#4a\draw[#5,line width=2.6pt,-{Latex[length=2.6mm]}] (\xa,0)--(\xb,0);\else\ifx#3a\draw[#5,line width=2.6pt,{Latex[length=2.6mm]}-] (\xa,0)--(\xb,0);\else\draw[#5,line width=2.6pt] (\xa,0)--(\xb,0);\fi\fi
  \intvdot{#1}{#3}{#5}\intvdot{#2}{#4}{#5}
  \ifx#3a\else\node[below=2pt,font=\scriptsize] at (#1,0) {$#1$};\fi
  \ifx#4a\else\node[below=2pt,font=\scriptsize] at (#2,0) {$#2$};\fi
\end{tikzpicture}}

\newcommand{\passos}[1]{\begin{tcolorbox}[fmbase,colback=fmdefbg,colframe=fmdef,coltitle=white,title={\faListOl\ \ Com es resolen}]#1\end{tcolorbox}}
\newcommand{\passosT}[2]{\begin{tcolorbox}[fmbase,colback=fmdefbg,colframe=fmdef,coltitle=white,title={\faListOl\ \ #1}]#2\end{tcolorbox}}
% núvol: eixos amb graella; #1 xmax, #2 ymax, #3 nom X, #4 nom Y
\newcommand{\nuvol}[4]{\draw[fmmuted!25] (0,0) grid (#1,#2);
  \draw[-{Latex[length=2mm]},fmmuted,line width=0.8pt] (0,0) -- (#1+0.6,0) node[right,font=\scriptsize] {#3};
  \draw[-{Latex[length=2mm]},fmmuted,line width=0.8pt] (0,0) -- (0,#2+0.6) node[above,font=\scriptsize] {#4};}
\newcommand{\pnt}[1]{\fill[fmblue] #1 circle (3pt);}


\begin{document}
\portada{Dues variables alhora: estan relacionades? Podem preveure l'una a partir de l'altra?}

{\color{fmnavy}\heavy\Large\faMapSigns\ \ Què aprendràs?}\par\smallskip
\begin{minipage}[t]{0.485\linewidth}
\itile{sec1}{10.1}{\faTable}{Distribucions bidimensionals}
\itile{sec2}{10.2}{\faBraille}{Núvol de punts}
\itile{sec3}{10.3}{\faChartLine}{Correlació}
\itile{sec4}{10.4}{\faRulerHorizontal}{Coeficient de correlació}
\end{minipage}\hfill
\begin{minipage}[t]{0.485\linewidth}
\itile{sec5}{10.5}{\faRoute}{Recta de regressió}
\itile{sec6}{10.6}{\faBullseye}{Estimacions}
\itile{resum}{\phantom{10.6}}{\faClipboardList}{En resum}
\end{minipage}
\llegenda

% ============================================================ 10.1
\clearpage
\seccio{sec1}{Distribucions bidimensionals}{Dues característiques de cada individu: estan relacionades?}

\vfill
\begin{idea}[Variable bidimensional $(X,Y)$]
Estudia dues variables $X$ i $Y$ sobre els \textbf{mateixos individus}. Cada individu dona una \textbf{parella} $(x_i,\,y_i)$; el conjunt de totes és una \textbf{distribució bidimensional}, amb $N$ parelles.
\end{idea}

\vfill
\begin{exemple}[Exemple: \ hores d'estudi ($X$) i nota ($Y$) de 10 alumnes]
{\centering\small \taularod[fmex]{>{\heavy}l*{10}{>{\hspace{3pt}}c<{\hspace{3pt}}}}{
\rowcolor{fmex}\color{white}Hores $x_i$ & \color{white}1 & \color{white}2 & \color{white}2 & \color{white}3 & \color{white}3 & \color{white}4 & \color{white}5 & \color{white}5 & \color{white}6 & \color{white}7\\
Nota $y_i$ & 3 & 3 & 5 & 4 & 6 & 6 & 7 & 8 & 8 & 9
}\par}\medskip
La primera alumna ha estudiat 1 hora i ha tret un 3: la parella $(1,\,3)$. \ $N=10$ parelles.
\end{exemple}

\vfill
\begin{compte}[L'ordre importa]
$(2,\,5)$ vol dir «2 hores i un 5», que no és el mateix que $(5,\,2)$, «5 hores i un 2».
\end{compte}

\vfill
\begin{exemple}[Exemple: \ fills ($X$) i cotxes ($Y$) de 20 famílies: \ taula de doble entrada]
\begin{minipage}[c]{0.42\linewidth}\centering\small
\taularod[fmex]{>{\heavy}c*{4}{>{\hspace{6pt}}c<{\hspace{6pt}}}}{
\rowcolor{fmex}\color{white}$X\backslash Y$ & \color{white}0 & \color{white}1 & \color{white}2 & \color{white}$f_x$\\
1 & 2 & 4 & 1 & \textbf{7}\\
\rowcolor{fmexbg}2 & 1 & \cellcolor{fmyellow!45}5 & 3 & \textbf{9}\\
3 & 0 & 1 & 3 & \textbf{4}\\
\rowcolor{fmexbg}$f_y$ & \textbf{3} & \textbf{10} & \textbf{7} & \textbf{20}
}
\end{minipage}\hfill
\begin{minipage}[c]{0.55\linewidth}
Si les parelles es repeteixen: $X$ en les files, $Y$ en les columnes i, en cada casella, la \textbf{freqüència} de la parella.\\[3pt]
5 famílies tenen 2 fills i 1 cotxe.\\[3pt]
Columna $f_x$: la distribució de $X$ sola. Fila $f_y$: la de $Y$. Totes dues sumen $N=20$.
\end{minipage}
\end{exemple}

\vfill
\begin{tcbraster}[raster columns=3,raster equal height,raster column skip=8pt]
\begin{tcolorbox}[fmbase,colback=fmK!8,colframe=fmK,coltitle=white,title={Relació funcional}]
Hi ha una \textbf{fórmula}: $X$ determina $Y$.\par\smallskip
\small Costat i perímetre d'un quadrat: $y=4x$.
\end{tcolorbox}
\begin{tcolorbox}[fmbase,colback=fmA!8,colframe=fmA,coltitle=white,title={Relació estadística}]
Hi ha una \textbf{tendència}, però no una fórmula.\par\smallskip
\small Alçada i pes: els més alts solen pesar més.
\end{tcolorbox}
\begin{tcolorbox}[fmbase,colback=fmR!6,colframe=fmR,coltitle=white,title={Cap relació}]
Saber $X$ no diu res de $Y$: \textbf{independents}.\par\smallskip
\small Número de calçat i nota de matemàtiques.
\end{tcolorbox}
\end{tcbraster}

\vfill
% ============================================================ 10.2
\clearpage
\seccio{sec2}{Núvol de punts}{Cada parella, un punt del pla.}

\vfill
\begin{minipage}[c]{0.52\linewidth}
\passosT{Com es dibuixa}{\pas{1} Una \textbf{escala} per a cada eix que arreplegue totes les dades (no cal que siga la mateixa). El nom de cada variable en el seu eix.\\[3pt]
\pas{2} Un \textbf{punt} per a cada parella $(x_i,y_i)$.\\[3pt]
\pas{3} Si una parella es repeteix, escriu al costat quantes vegades (o fes el punt més gros).}
\end{minipage}\hfill
\begin{minipage}[c]{0.45\linewidth}\centering
\begin{tikzpicture}[x=0.55cm,y=0.42cm]
  \nuvol{8}{10}{hores}{nota}
  \foreach \k in {2,4,6,8} \node[below,font=\scriptsize] at (\k,0) {\k};
  \foreach \k in {2,4,...,10} \node[left,font=\scriptsize] at (0,\k) {\k};
  \foreach \p in {(1,3),(2,3),(2,5),(3,4),(3,6),(4,6),(5,7),(5,8),(6,8),(7,9)} {\pnt{\p}}
  %
\end{tikzpicture}\\
{\comic\small\color{fmmuted}els punts «pugen» cap a la dreta}
\end{minipage}

\vfill
\begin{truc}[Dades grans]
Amb alçades de 150 a 190 cm, l'eix no ha de començar en 0: comença un poc abans de la dada més xicoteta i marca el tall de l'eix amb un \textbf{zigzag}.
\end{truc}

\vfill
\begin{minipage}[c]{0.5\linewidth}
\begin{idea}[Centre de gravetat]
El punt de les dues mitjanes, $G=(\bar x,\,\bar y)$, és el «punt mitjà» del núvol:\\[3pt]
{\centering$\bar x=\dfrac{\sum x_i}{N}\qquad\bar y=\dfrac{\sum y_i}{N}$\par}
\end{idea}
\end{minipage}\hfill
\begin{minipage}[c]{0.47\linewidth}\centering
\begin{tikzpicture}[x=0.55cm,y=0.42cm]
  \nuvol{8}{10}{hores}{nota}
  \foreach \k in {2,4,6,8} \node[below,font=\scriptsize] at (\k,0) {\k};
  \foreach \k in {2,4,...,10} \node[left,font=\scriptsize] at (0,\k) {\k};
  \foreach \p in {(1,3),(2,3),(2,5),(3,4),(3,6),(4,6),(5,7),(5,8),(6,8),(7,9)} {\pnt{\p}}
  \draw[fmR,dashed] (3.8,0) -- (3.8,10) (0,5.9) -- (8,5.9);
  \filldraw[fill=white,draw=fmR,line width=1.4pt] (3.8,5.9) circle (5pt) node[above right,font=\scriptsize\heavy,text=fmR] {$G$};
\end{tikzpicture}
\end{minipage}

\vfill
\begin{exemple}[Exemple: \ el centre de gravetat de les hores d'estudi i les notes]
$\bar x=\dfrac{1+2+2+3+3+4+5+5+6+7}{10}=\dfrac{38}{10}=3{,}8$ \qquad $\bar y=\dfrac{3+3+5+4+6+6+7+8+8+9}{10}=\dfrac{59}{10}=5{,}9$ \qquad $G=\boxed{(3{,}8,\ 5{,}9)}$\\[4pt]
Les línies $x=\bar x$ i $y=\bar y$ fan quatre zones: quasi tots els punts estan a baix a l'esquerra o a dalt a la dreta. Les variables «van juntes».
\end{exemple}

\vfill
\begin{exemple}[Exemple: \ $x_i$: 2, 4, 5, 6, 8 \ i \ $y_i$: 7, 6, 4, 3, 0]
$\bar x=\dfrac{25}{5}=5$, \ $\bar y=\dfrac{20}{5}=4$: \ $G=\boxed{(5,\,4)}$. \ Ara els punts «baixen»: quan $x$ creix, $y$ decreix.
\end{exemple}

\vfill
% ============================================================ 10.3
\clearpage
\seccio{sec3}{Correlació}{La forma del núvol diu com és la relació.}

\vfill
\begin{idea}[Correlació]
La \textbf{correlació} és la relació estadística entre dues variables.\\[3pt]
\textbf{Sentit}: \textbf{positiva} (directa) si, quan $X$ creix, $Y$ tendeix a créixer: el núvol puja. \textbf{Negativa} (inversa) si $Y$ tendeix a decréixer: el núvol baixa. \textbf{Nul·la} si no puja ni baixa.\\[3pt]
\textbf{Força}: \textbf{forta} si els punts estan molt a prop d'una recta; \textbf{feble} si estan escampats. Tots en una recta: relació \textbf{funcional}.
\end{idea}

\vfill
\begin{tcbraster}[raster columns=3,raster equal height,raster column skip=8pt,raster row skip=8pt]
\begin{tcolorbox}[fmbase,colback=fmK!8,colframe=fmK,coltitle=white,title={Positiva i forta},before upper=\centering,left=1mm,right=1mm]\begin{tikzpicture}[scale=0.32]\draw[fmmuted,line width=0.6pt] (0,0) rectangle (6,6);\fill[fmblue] (0.6,0.8) circle (3.5pt);\fill[fmblue] (1.2,1.4) circle (3.5pt);\fill[fmblue] (1.8,1.7) circle (3.5pt);\fill[fmblue] (2.4,2.6) circle (3.5pt);\fill[fmblue] (3,2.9) circle (3.5pt);\fill[fmblue] (3.6,3.7) circle (3.5pt);\fill[fmblue] (4.2,4.1) circle (3.5pt);\fill[fmblue] (4.8,4.9) circle (3.5pt);\fill[fmblue] (5.4,5.3) circle (3.5pt);\end{tikzpicture}\end{tcolorbox}
\begin{tcolorbox}[fmbase,colback=fmK!8,colframe=fmK,coltitle=white,title={Positiva i feble},before upper=\centering,left=1mm,right=1mm]\begin{tikzpicture}[scale=0.32]\draw[fmmuted,line width=0.6pt] (0,0) rectangle (6,6);\fill[fmblue] (0.6,1.6) circle (3.5pt);\fill[fmblue] (1.2,0.8) circle (3.5pt);\fill[fmblue] (1.8,2.9) circle (3.5pt);\fill[fmblue] (2.4,1.9) circle (3.5pt);\fill[fmblue] (3,3.8) circle (3.5pt);\fill[fmblue] (3.6,2.4) circle (3.5pt);\fill[fmblue] (4.2,4.6) circle (3.5pt);\fill[fmblue] (4.8,3.3) circle (3.5pt);\fill[fmblue] (5.4,5) circle (3.5pt);\end{tikzpicture}\end{tcolorbox}
\begin{tcolorbox}[fmbase,colback=fmmuted!8,colframe=fmmuted,coltitle=white,title={Nul·la},before upper=\centering,left=1mm,right=1mm]\begin{tikzpicture}[scale=0.32]\draw[fmmuted,line width=0.6pt] (0,0) rectangle (6,6);\fill[fmblue] (0.6,3.8) circle (3.5pt);\fill[fmblue] (1.2,1.2) circle (3.5pt);\fill[fmblue] (1.8,5) circle (3.5pt);\fill[fmblue] (2.4,2.6) circle (3.5pt);\fill[fmblue] (3,4.4) circle (3.5pt);\fill[fmblue] (3.6,0.9) circle (3.5pt);\fill[fmblue] (4.2,3.1) circle (3.5pt);\fill[fmblue] (4.8,5.3) circle (3.5pt);\fill[fmblue] (5.4,1.8) circle (3.5pt);\end{tikzpicture}\end{tcolorbox}
\begin{tcolorbox}[fmbase,colback=fmR!8,colframe=fmR,coltitle=white,title={Negativa i feble},before upper=\centering,left=1mm,right=1mm]\begin{tikzpicture}[scale=0.32]\draw[fmmuted,line width=0.6pt] (0,0) rectangle (6,6);\fill[fmblue] (0.6,4.4) circle (3.5pt);\fill[fmblue] (1.2,5.2) circle (3.5pt);\fill[fmblue] (1.8,3.1) circle (3.5pt);\fill[fmblue] (2.4,4.1) circle (3.5pt);\fill[fmblue] (3,2.2) circle (3.5pt);\fill[fmblue] (3.6,3.6) circle (3.5pt);\fill[fmblue] (4.2,1.4) circle (3.5pt);\fill[fmblue] (4.8,2.7) circle (3.5pt);\fill[fmblue] (5.4,1.0) circle (3.5pt);\end{tikzpicture}\end{tcolorbox}
\begin{tcolorbox}[fmbase,colback=fmR!8,colframe=fmR,coltitle=white,title={Negativa i forta},before upper=\centering,left=1mm,right=1mm]\begin{tikzpicture}[scale=0.32]\draw[fmmuted,line width=0.6pt] (0,0) rectangle (6,6);\fill[fmblue] (0.6,5.2) circle (3.5pt);\fill[fmblue] (1.2,4.6) circle (3.5pt);\fill[fmblue] (1.8,4.3) circle (3.5pt);\fill[fmblue] (2.4,3.4) circle (3.5pt);\fill[fmblue] (3,3.1) circle (3.5pt);\fill[fmblue] (3.6,2.3) circle (3.5pt);\fill[fmblue] (4.2,1.9) circle (3.5pt);\fill[fmblue] (4.8,1.1) circle (3.5pt);\fill[fmblue] (5.4,0.7) circle (3.5pt);\end{tikzpicture}\end{tcolorbox}
\begin{tcolorbox}[fmbase,colback=fmA!8,colframe=fmA,coltitle=white,title={Funcional},before upper=\centering,left=1mm,right=1mm]\begin{tikzpicture}[scale=0.32]\draw[fmmuted,line width=0.6pt] (0,0) rectangle (6,6);\draw[fmR!70,line width=0.8pt] (0.3,0.3)--(5.6,5.6);\fill[fmblue] (0.6,0.6) circle (3.5pt);\fill[fmblue] (1.5,1.5) circle (3.5pt);\fill[fmblue] (2.4,2.4) circle (3.5pt);\fill[fmblue] (3.3,3.3) circle (3.5pt);\fill[fmblue] (4.2,4.2) circle (3.5pt);\fill[fmblue] (5.1,5.1) circle (3.5pt);\end{tikzpicture}\end{tcolorbox}
\end{tcbraster}

\vfill
\begin{minipage}[c]{0.55\linewidth}\centering\small
\taularod[fmex]{>{\hspace{3pt}}l<{\hspace{3pt}}>{\hspace{3pt}}l<{\hspace{3pt}}}{
\rowcolor{fmex}\color{white}\heavy Variables & \color{white}\heavy Correlació\\
Hores d'estudi i nota & positiva i forta\\
\rowcolor{fmexbg}Temperatura i gelats venuts & positiva i forta\\
Km d'un cotxe de segona mà i preu & negativa i forta\\
\rowcolor{fmexbg}Hores de pantalla i hores de son & negativa i feble\\
Alçada i nota de matemàtiques & nul·la
}
\end{minipage}\hfill
\begin{minipage}[c]{0.42\linewidth}
\begin{truc}[Lineal o no lineal]
Si els punts segueixen una \textbf{corba} (el rendiment puja amb l'entrenament fins que es cansa i baixa), la relació pot ser forta però \textbf{no lineal}. Ací estudiem les lineals.
\end{truc}
\end{minipage}

\vfill
\begin{compte}[Correlació no vol dir causa]
A l'estiu es venen més gelats i també s'ofeguen més persones a la platja: correlació positiva, però els gelats no en són la causa. Totes dues depenen d'una tercera variable: la \textbf{calor}. Pregunta't si hi ha una altra variable que les faça variar totes dues.
\end{compte}

\vfill
% ============================================================ 10.4
\clearpage
\seccio{sec4}{Coeficient de correlació}{Un nombre entre $-1$ i $1$ que mesura la correlació.}

\vfill
\begin{minipage}[c]{0.45\linewidth}
\begin{idea}[El coeficient $r$ (de Pearson)]
Sempre $-1\le r\le1$.\\[3pt]
El \textbf{signe} diu el sentit: $r>0$ positiva; $r<0$ negativa.\\[3pt]
$|r|$ diu la \textbf{força}: a prop d'1, forta; a prop de 0, feble.
\end{idea}
\end{minipage}\hfill
\begin{minipage}[c]{0.52\linewidth}\centering
\begin{tikzpicture}[x=0.6cm]
  \shade[left color=fmR!60,right color=fmK!60] (-5,0) rectangle (5,0.35);
  \fill[white,opacity=0.75] (-2.5,0) rectangle (2.5,0.35);
  \draw[fmnavy,line width=0.8pt] (-5,0) rectangle (5,0.35);
  \foreach \x/\l in {-5/-1,-4/{-0,8},-2.5/{-0,5},0/0,2.5/{0,5},4/{0,8},5/1} {\draw[fmnavy] (\x,0) -- (\x,-0.12); \node[below,font=\scriptsize] at (\x,-0.12) {$\l$};}
  \node[font=\scriptsize,align=center] at (-4.5,0.8) {forta}; \node[font=\scriptsize] at (-1.25,0.8) {feble}; \node[font=\scriptsize] at (1.25,0.8) {feble}; \node[font=\scriptsize] at (4.5,0.8) {forta};
  \node[font=\scriptsize\heavy,text=fmR] at (-3.5,-0.9) {negativa}; \node[font=\scriptsize\heavy,text=fmK!80!black] at (3.5,-0.9) {positiva};
\end{tikzpicture}\\[4pt]
{\small $|r|\ge0{,}8$ forta \ \ $0{,}5\le|r|<0{,}8$ moderada \ \ $|r|<0{,}5$ feble}
\end{minipage}

\vfill
\begin{compte}[La força la dona $|r|$]
$r=-0{,}9$ és una correlació \textbf{més forta} que $r=0{,}6$: el signe només diu si el núvol puja o baixa.
\end{compte}

\vfill
\begin{minipage}[c]{0.5\linewidth}
\passosT{Com es calcula}{\pas{1} Taula amb $x_i$, $y_i$, $x_i^2$, $y_i^2$, $x_iy_i$ i les sumes.\\[3pt]
\pas{2} Mitjanes: \ $\bar x=\frac{\sum x_i}{N}$, \ $\bar y=\frac{\sum y_i}{N}$.\\[3pt]
\pas{3} $\sigma_x=\sqrt{\frac{\sum x_i^2}{N}-\bar x^2}$, \ $\sigma_y=\sqrt{\frac{\sum y_i^2}{N}-\bar y^2}$.\\[3pt]
\pas{4} Covariància: \ $\sigma_{xy}=\frac{\sum x_iy_i}{N}-\bar x\,\bar y$.}
\end{minipage}\hfill
\begin{minipage}[c]{0.47\linewidth}
\begin{tcolorbox}[enhanced,colback=fmyellow!18,colframe=fmgold,arc=4mm,boxrule=1.6pt,before upper=\centering,top=2mm,bottom=2mm]
{\Large$r=\dfrac{\sigma_{xy}}{\sigma_x\cdot\sigma_y}$}\par\smallskip
\small $\sigma_{xy}>0$: creixen juntes; \ $\sigma_{xy}<0$: l'una creix i l'altra decreix
\end{tcolorbox}
\end{minipage}

\vfill
\begin{exemple}[Exemple: \ hores d'estudi i nota de 5 alumnes]
\begin{minipage}[c]{0.44\linewidth}\centering\small
\taularod[fmex]{*{5}{>{\hspace{3pt}}c<{\hspace{3pt}}}}{
\rowcolor{fmex}\color{white}$x_i$ & \color{white}$y_i$ & \color{white}$x_i^2$ & \color{white}$y_i^2$ & \color{white}$x_iy_i$\\
1 & 2 & 1 & 4 & 2\\
\rowcolor{fmexbg}2 & 3 & 4 & 9 & 6\\
3 & 5 & 9 & 25 & 15\\
\rowcolor{fmexbg}4 & 4 & 16 & 16 & 16\\
5 & 6 & 25 & 36 & 30\\
\rowcolor{fmexbg}\textbf{15} & \textbf{20} & \textbf{55} & \textbf{90} & \textbf{69}
}
\end{minipage}\hfill
\begin{minipage}[c]{0.53\linewidth}
$\bar x=3$, \ $\bar y=4$\\[3pt]
$\sigma_x=\sqrt{\frac{55}5-9}=\sqrt2$, \ $\sigma_y=\sqrt{\frac{90}5-16}=\sqrt2$\\[3pt]
$\sigma_{xy}=\frac{69}5-3\cdot4=13{,}8-12=1{,}8$\\[3pt]
$r=\dfrac{1{,}8}{\sqrt2\cdot\sqrt2}=\boxed{0{,}9}$: \ positiva i \textbf{forta}
\end{minipage}
\end{exemple}

\vfill
\begin{exemple}[Exemple: \ preu (€) i unitats venudes en una setmana]
Preus 1, 2, 3, 4, 5 \ i \ unitats 50, 44, 40, 32, 24: \ $\sum y_i=190$, \ $\sum y_i^2=7636$, \ $\sum x_iy_i=506$\\[4pt]
$\bar y=38$, \ $\sigma_y=\sqrt{\frac{7636}5-38^2}=\sqrt{83{,}2}\approx9{,}12$ \qquad $\sigma_{xy}=\frac{506}5-3\cdot38=-12{,}8$ \qquad $r=\dfrac{-12{,}8}{\sqrt2\cdot9{,}12}\approx\boxed{-0{,}99}$\\[3pt]
Negativa i \textbf{molt forta}: com més car, menys se'n ven.
\end{exemple}

\vfill
\begin{exemple}[Exemple: \ les hores d'estudi i les notes dels 10 alumnes (10.1)]
\begin{minipage}[c]{0.5\linewidth}
$\sum x_i=38$, \ $\sum y_i=59$, \ $\sum x_i^2=178$, \ $\sum y_i^2=389$, \ $\sum x_iy_i=259$\\[4pt]
$\bar x=3{,}8$, \ $\bar y=5{,}9$\\[3pt]
$\sigma_x=\sqrt{17{,}8-3{,}8^2}=\sqrt{3{,}36}\approx1{,}83$\\[3pt]
$\sigma_y=\sqrt{38{,}9-5{,}9^2}=\sqrt{4{,}09}\approx2{,}02$
\end{minipage}\hfill
\begin{minipage}[c]{0.47\linewidth}
$\sigma_{xy}=25{,}9-3{,}8\cdot5{,}9=3{,}48$\\[4pt]
$r=\dfrac{3{,}48}{1{,}83\cdot2{,}02}\approx\boxed{0{,}94}$\\[4pt]
Positiva i forta, com dèiem en mirar el núvol (10.3).
\end{minipage}
\end{exemple}

\vfill
\begin{exemple}[Exemple: \ $x_i$: 1, 2, 3, 4, 5 \ i \ $y_i$: 3, 1, 4, 2, 5]
\begin{minipage}[c]{0.3\linewidth}\centering
\begin{tikzpicture}[x=0.5cm,y=0.42cm]
  \nuvol{6}{6}{$x$}{$y$}
  \foreach \p in {(1,3),(2,1),(3,4),(4,2),(5,5)} {\pnt{\p}}
\end{tikzpicture}
\end{minipage}\hfill
\begin{minipage}[c]{0.67\linewidth}
$\bar x=3$, \ $\bar y=3$, \ $\sigma_x=\sigma_y=\sqrt2$ \qquad $\sum x_iy_i=3+2+12+8+25=50$\\[4pt]
$\sigma_{xy}=\frac{50}5-3\cdot3=1$ \qquad $r=\dfrac{1}{\sqrt2\cdot\sqrt2}=\boxed{0{,}5}$\\[4pt]
Positiva però només \textbf{moderada}: el núvol puja, però els punts estan escampats.
\end{minipage}
\end{exemple}

\vfill
\begin{tcbraster}[raster columns=2,raster equal height,raster column skip=10pt]
\begin{truc}[Amb la calculadora]
En el mode estadístic de dues variables ($y=a+bx$) introdueix les parelles: dona $\bar x$, $\bar y$, $\sigma_x$, $\sigma_y$ i $r$. Fes-la servir per a comprovar.
\end{truc}
\begin{compte}[Compte amb l'ordre]
En $\sigma_{xy}$ primer dividim $\sum x_iy_i$ entre $N$ i després restem $\bar x\,\bar y$. Si $|r|$ et dona més gran que 1, t'has equivocat.
\end{compte}
\end{tcbraster}

\vfill
% ============================================================ 10.5
\clearpage
\seccio{sec5}{Recta de regressió}{La recta que millor s'ajusta al núvol.}

\vfill
\begin{minipage}[c]{0.55\linewidth}
\begin{idea}[Recta de regressió de $Y$ sobre $X$]
Si la correlació és forta, els punts s'agrupen al voltant d'una recta. La que millor s'hi ajusta resumeix la relació i serveix per a fer \textbf{estimacions}.\\[3pt]
\textbf{A ull}: per $G=(\bar x,\bar y)$, seguint el núvol, amb tants punts per damunt com per davall.\\[3pt]
\textbf{Amb la fórmula}: la que fa mínims els errors (al quadrat).
\end{idea}
\end{minipage}\hfill
\begin{minipage}[c]{0.42\linewidth}
\begin{tcolorbox}[enhanced,colback=fmyellow!18,colframe=fmgold,arc=4mm,boxrule=1.6pt,before upper=\centering,top=2mm,bottom=2mm]
{\Large$y-\bar y=\dfrac{\sigma_{xy}}{\sigma_x^2}\,(x-\bar x)$}\par\smallskip
\small punt-pendent (unitat 7): passa per $G$ amb pendent $m=\dfrac{\sigma_{xy}}{\sigma_x^2}$
\end{tcolorbox}
\end{minipage}

\vfill
\begin{truc}[El signe del pendent]
$m$ té el mateix signe que la covariància i que $r$: correlació positiva, la recta \textbf{puja}; negativa, \textbf{baixa}. La calculadora (mode $y=a+bx$) dona $a=n$ i $b=m$.
\end{truc}

\vfill
\begin{exemple}[Exemple: \ hores d'estudi i nota de 5 alumnes (10.4)]
\begin{minipage}[c]{0.36\linewidth}\centering
\begin{tikzpicture}[x=0.55cm,y=0.42cm]
  \nuvol{6}{7}{hores}{nota}
  \foreach \k in {1,...,5} \node[below,font=\scriptsize] at (\k,0) {\k};
  \foreach \k in {2,4,6} \node[left,font=\scriptsize] at (0,\k) {\k};
  \draw[fmR,line width=1.4pt] (0,1.3) -- (6,6.7);
  \foreach \p in {(1,2),(2,3),(3,5),(4,4),(5,6)} {\pnt{\p}}
  \filldraw[fill=white,draw=fmR,line width=1.2pt] (3,4) circle (4pt);
\end{tikzpicture}
\end{minipage}\hfill
\begin{minipage}[c]{0.62\linewidth}
$\bar x=3$, \ $\bar y=4$, \ $\sigma_x^2=2$, \ $\sigma_{xy}=1{,}8$: \ $m=\dfrac{1{,}8}{2}=0{,}9$\\[4pt]
$y-4=0{,}9\,(x-3)\ \to\ \boxed{y=0{,}9x+1{,}3}$\\[4pt]
Per cada hora més d'estudi, la nota puja, de mitjana, \textbf{0,9 punts}.
\end{minipage}
\end{exemple}

\vfill
\begin{exemple}[Exemple: \ preu i unitats venudes (10.4)]
$\bar x=3$, \ $\bar y=38$, \ $\sigma_x^2=2$, \ $\sigma_{xy}=-12{,}8$: \ $m=\dfrac{-12{,}8}{2}=-6{,}4$\\[4pt]
$y-38=-6{,}4\,(x-3)\ \to\ \boxed{y=-6{,}4x+57{,}2}$: \ per cada euro més, es venen unes \textbf{6,4 unitats menys}.
\end{exemple}

\vfill
\begin{exemple}[Exemple: \ $x_i$: 2, 4, 5, 6, 8 \ i \ $y_i$: 7, 6, 4, 3, 0 \ (10.2)]
$\sum x_i^2=145$, \ $\sum x_iy_i=81$: \ $\sigma_x^2=\frac{145}5-25=4$, \ $\sigma_{xy}=\frac{81}5-5\cdot4=-3{,}8$\\[4pt]
$m=\dfrac{-3{,}8}{4}=-0{,}95$: \ $y-4=-0{,}95\,(x-5)\ \to\ \boxed{y=-0{,}95x+8{,}75}$ \ \nota{recta que baixa}
\end{exemple}

\vfill
% ============================================================ 10.6
\clearpage
\seccio{sec6}{Estimacions}{Preveure $Y$ per a un valor de $X$ que no tenim.}

\vfill
\begin{minipage}[c]{0.52\linewidth}
\begin{idea}[Estimar amb la recta]
Amb $y=mx+n$, per a un valor $x_0$:\\[3pt]
{\centering\Large$\hat y=m\cdot x_0+n$\par}\smallskip
$\hat y$ és una \textbf{estimació}: el valor que esperem de mitjana, no el que tindrà exactament un individu.
\end{idea}
\end{minipage}\hfill
\begin{minipage}[c]{0.45\linewidth}
\passosT{Quan és fiable?}{\pas{1} Correlació \textbf{forta}: $|r|\ge0{,}8$, bones; entre 0,5 i 0,8, aproximades; menys, no fiables.\\[3pt]
\pas{2} $x_0$ \textbf{dins} de les dades: \textbf{interpolar}. Fora: \textbf{extrapolar}, menys fiable.\\[3pt]
\pas{3} \textbf{Moltes} dades.}
\end{minipage}

\vfill
\begin{exemple}[Exemple: \ hores d'estudi i nota, \ $y=0{,}9x+1{,}3$ \ ($r=0{,}9$)]
3,5 hores: \ $\hat y=0{,}9\cdot3{,}5+1{,}3=\boxed{4{,}45}$ \ \nota{fiable: $r$ alt i 3,5 dins d'1 a 5 hores}\\[4pt]
12 hores: \ $\hat y=0{,}9\cdot12+1{,}3=12{,}1$: \ \mal{\text{una nota impossible}} \ Les dades van d'1 a 5 hores: no sabem si tan lluny la relació continua igual.
\end{exemple}

\vfill
\begin{exemple}[Exemple: \ preu i vendes, \ $y=-6{,}4x+57{,}2$ \ ($r\approx-0{,}99$)]
\begin{minipage}[c]{0.6\linewidth}
3,50 €: \ $\hat y=-6{,}4\cdot3{,}5+57{,}2=34{,}8\approx\boxed{35\text{ unitats}}$ \ \nota{molt fiable}\\[4pt]
10 €: \ $-6{,}4\cdot10+57{,}2=-6{,}8$ unitats: \ impossible, extrapolem massa.\\[4pt]
\textbf{Estimar $x$}: per a vendre'n 45, \ $45=-6{,}4x+57{,}2$:\\ $x=\dfrac{12{,}2}{6{,}4}\approx\boxed{1{,}91\text{ €}}$ \ \nota{només si $|r|$ és quasi 1}
\end{minipage}\hfill
\begin{minipage}[c]{0.37\linewidth}\centering
\begin{tikzpicture}[x=0.5cm,y=0.05cm]
  \draw[fmmuted!25] (0,0) grid[xstep=1,ystep=10] (6,60);
  \draw[-{Latex[length=2mm]},fmmuted,line width=0.8pt] (0,0) -- (6.6,0) node[right,font=\scriptsize] {€};
  \draw[-{Latex[length=2mm]},fmmuted,line width=0.8pt] (0,0) -- (0,65) node[above,font=\scriptsize] {unitats};
  \foreach \k in {1,...,5} \node[below,font=\scriptsize] at (\k,0) {\k};
  \foreach \k in {20,40,60} \node[left,font=\scriptsize] at (0,\k) {\k};
  \draw[fmR,line width=1.3pt] (0.3,55.3) -- (6,18.8);
  \foreach \p in {(1,50),(2,44),(3,40),(4,32),(5,24)} {\pnt{\p}}
  \draw[fmK,dashed] (3.5,0) -- (3.5,34.8) -- (0,34.8); \fill[fmK] (3.5,34.8) circle (3pt);
\end{tikzpicture}
\end{minipage}
\end{exemple}

\vfill
\begin{exemple}[Exemples]
Recta $y=-0{,}95x+8{,}75$ (10.5): \ per a $x=3$, \ $\hat y=5{,}9$; \ per a $x=7$, \ $\hat y=2{,}1$ \ \nota{interpolacions}\\[4pt]
Si $r=0{,}35$, cap estimació és fiable: els punts estan massa escampats al voltant de la recta.
\end{exemple}

\vfill
\begin{compte}[Interpolar sí, extrapolar amb compte]
Lluny de les dades, la recta pot donar resultats sense sentit (notes de més de 10, vendes negatives). Comprova sempre que l'estimació és raonable.
\end{compte}

\vfill
% ============================================================ resum
\resumtitol
\begin{tcbraster}[raster columns=2,raster equal height=rows,raster column skip=14pt,raster row skip=40pt,raster before skip=20pt]
\rtile{10.1}{Bidimensional}{Parelles $(x_i,y_i)$. Taula simple o de doble entrada.}
\rtile{10.2}{Núvol de punts}{Un punt per parella. Centre de gravetat $G=(\bar x,\bar y)$.}
\rtile{10.3}{Correlació}{Positiva, negativa o nul·la; forta o feble. No vol dir causa.}
\rtile{10.4}{Coeficient}{$r=\frac{\sigma_{xy}}{\sigma_x\sigma_y}$, \ $-1\le r\le1$. \ $\sigma_{xy}=\frac{\sum x_iy_i}{N}-\bar x\bar y$.}
\rtile{10.5}{Regressió}{$y-\bar y=\frac{\sigma_{xy}}{\sigma_x^2}(x-\bar x)$.}
\rtile{10.6}{Estimacions}{$\hat y=mx_0+n$. Fiables si $|r|$ és gran i $x_0$ és dins de les dades.}
\end{tcbraster}
\end{document}
